\section{One-sided criteria for zeros}
\label{sec:one-sided}

\subsection{A shorter sufficient criterion for the Riemann hypothesis}

For the zeta function alone, an eventual bound admits a shorter proof.
The hypothesis here is stronger than a bound on arbitrarily late
intervals; it has the advantage that either direction of inequality
is sufficient.

\begin{lemma}[The real boundary for a positive Mellin transform]
\label{rh-positive-mellin-boundary}
Let $X\ge1$, and let $f:[X,\infty)\to[0,\infty)$ be locally integrable
and of at most polynomial growth. If its Mellin transform
\[
 F(s)=\int_X^\infty f(x)x^{-s-1}\,dx
\]
has a finite abscissa of convergence $\alpha$, it does not extend
holomorphically through the real point $\alpha$.
\end{lemma}
\begin{proof}
Suppose that it extends to a disk of radius $r>0$ about $\alpha$,
and put $\sigma_0=\alpha+r/4$. Its Taylor series at $\sigma_0$,
evaluated at $\sigma_0-z$, converges for $r/4<z<r/2$ and has terms
\[
 \frac{(-1)^kF^{(k)}(\sigma_0)z^k}{k!}
 =\int_X^\infty f(x)x^{-\sigma_0-1}
                   \frac{(z\log x)^k}{k!}\,dx\ge0.
\]
Differentiation follows by domination using convergence at an
intermediate real point between $\alpha$ and $\sigma_0$.
Tonelli's theorem identifies the Taylor sum with
$\int_X^\infty f(x)x^{-(\sigma_0-z)-1}\,dx$.
This is finite to the left of $\alpha$, a contradiction.
\end{proof}

\begin{proposition}[Either eventual one-sided zeta bound]
\label{rh-eventual-one-sided}
Fix $0<\theta<1$. If, for some finite $C\ge0$ and $X\ge1$, either
\[
 G_1(x)\le Cx^\theta\quad(x\ge X),
 \qquad\hbox{or}\qquad
 G_1(x)\ge-Cx^\theta\quad(x\ge X),
\]
then $\zeta(s)$ has no zero with $\Re s>\theta$.
Its pole at one is excluded from the zero-free assertion.
\end{proposition}
\begin{proof}
Let $\mathcal G(s)$ be the transform in
\eqref{eq:mellin-zeta}. It is holomorphic at one. It is also
holomorphic near every other positive real point: for $0<s<1$,
\[
 \sum_{j\ge0}\bigl((2j+1)^{-s}-(2j+2)^{-s}\bigr)>0,
 \qquad
 \zeta(s)=\frac{\sum_{n\ge1}(-1)^{n-1}n^{-s}}{1-2^{1-s}}<0,
\]
and the Euler product handles $s>1$.

In the upper-bound case put $f=Cx^\theta-G_1(x)$ on $[X,\infty)$;
in the lower-bound case put $f=Cx^\theta+G_1(x)$.
In either case $f\ge0$ and its Mellin transform, initially in a
right half-plane, is
\[
 \frac{CX^{\theta-s}}{s-\theta}\mp\mathcal G(s)
 \quad\hbox{plus an entire finite-interval integral}.
\]
This expression is holomorphic near every real $s>\theta$.
Lemma~\ref{rh-positive-mellin-boundary} therefore puts its abscissa
of convergence at most $\theta$; the case of abscissa $-\infty$
already has this conclusion. The integral is holomorphic on
$\Re s>\theta$, and hence so is $\mathcal G$.
But a zeta zero $\rho$ there would give the nonzero residue
$-m(\rho)/(\rho(\rho-1))$ in Lemma~\ref{lem:mellin}.
This contradiction proves the assertion.
\end{proof}

\begin{corollary}\label{rh-eventual-vanishing-loss}
Suppose $\delta_j>0$ tends to zero and $1/2+\delta_j<1$.
For every $j$, assume one of the two eventual bounds in
Proposition~\ref{rh-eventual-one-sided} with
$\theta=1/2+\delta_j$. The finite constant, threshold and direction
may depend on $j$. Then the Riemann hypothesis follows.
A single eventual bound with exponent exactly $1/2$ also suffices.
\end{corollary}
\begin{proof}
The proposition excludes every zero to the right of the critical line.
The functional equation excludes every nontrivial zero to its left.
\end{proof}

This argument does not replace an eventual bound by a cofinal-window
bound: nonnegativity on the entire tail was used in the Mellin integral.
For the latter arithmetic hypothesis one retains
Theorem~\ref{thm:vanishing-loss-windows}.

\subsection{Positive families on large intervals}

Only elementary complex analysis is needed in the following argument.
The proof is included because replacing a one-sided estimate by an
unjustified absolute-value estimate would invalidate the reduction.


\begin{proposition}\label{prop:landau-independent}
Let $g:[0,\infty)\to\R$ be measurable and satisfy
$|g(u)|\le C_0e^{2u}$. Suppose that its Laplace transform
\[
 \Phi(s)=\int_0^\infty g(u)e^{-su}\,du\qquad(\Rea s>2)
\]
extends meromorphically to $\Rea s>0$, with only isolated simple poles,
and that every real pole has positive real residue.
Fix $0<\theta<1$ and $\delta>0$. Suppose that for every $\eta>0$
there are arbitrarily large $U$ and a nonnegative function $R_U\ge g$
on $[U,(1+\delta)U]$ such that $R_U(u)e^{-\theta u}$ is integrable there and
\[
 \int_U^{(1+\delta)U}R_U(u)e^{-\theta u}\,du\le e^{\eta U}.
\]
Then $\Phi$ has no pole with real part greater than $\theta$.
\end{proposition}

\begin{proof}
We give the tail estimate and the pole argument explicitly.

Put $a=1+\delta/2$ and $b=1+\delta$. Both numbers
\[
 q_- = \frac{\exp(1-1/a)}{a},\qquad
 q_+ = \frac ba\exp(1-b/a)
\]
are strictly less than $1$. For the first, use
$\log a>1-1/a$ for $a>1$; for the second, use
$\log r<r-1$ for $r=b/a>1$. Choose $\lambda>0$ so large that
\[
 \lambda a\ge1,\qquad
 q:=\max(q_-,q_+)<\frac{\lambda}{\lambda+2},
 \qquad \sigma:=\lambda+2>2.
\]
For $U\ge1$, set $k=\lceil\lambda aU\rceil$, so that $k\ge U$.
Elementary exponential tilting and
$\int_0^\infty u^ke^{-tu}\,du=k!t^{-k-1}$ give
\begin{align*}
 \int_0^U\frac{u^ke^{-\lambda u}}{k!}\,du
 &\le \frac{e^{\lambda(a-1)U}}{(\lambda a)^{k+1}}
 \le \frac1{\lambda a}\left(\frac{q_-}{\lambda}\right)^k,\\
 \int_{bU}^\infty\frac{u^ke^{-\lambda u}}{k!}\,du
 &\le \frac{e^{-\lambda(b-a)U}}{(\lambda a/b)^{k+1}}
 \le \frac{b\exp((b-a)/a)}{\lambda a}
          \left(\frac{q_+}{\lambda}\right)^k.
\end{align*}
For the first inequality, write the integrand as
$e^{\lambda(a-1)u}u^ke^{-\lambda au}/k!$ and bound the first factor
by its value at $U$. For the second, write it as
$e^{-\lambda(1-a/b)u}u^ke^{-(\lambda a/b)u}/k!$ and bound the first
factor by its value at $bU$. The final bounds use
$U\le k/(\lambda a)$ and $U\ge(k-1)/(\lambda a)$, respectively.
Thus, for a constant $C_1$ depending only on the chosen $\delta,\lambda$,
\begin{equation}\label{eq:landau-tail-independent}
 \int_{[0,\infty)\setminus[U,bU]}
      \frac{u^ke^{-(\sigma-2)u}}{k!}\,du
 \le C_1\sigma^{-k}.
\end{equation}

For $\Rea s>2$, differentiation under the integral is justified by
the exponential bound on $g$. Write
\[
 A_k(s)=\frac{(-1)^k}{k!}\Phi^{(k)}(s)
       =\int_0^\infty\frac{g(u)u^ke^{-su}}{k!}\,du.
\]
The triangle inequality gives, for every real $\gamma$,
\[
 |A_k(\sigma+i\gamma)|+\Rea A_k(\sigma)
 \le\int_0^\infty (|g(u)|+g(u))
               \frac{u^ke^{-\sigma u}}{k!}\,du.
\]
On $[U,bU]$, use $|g|+g\le2R_U$ and
\[
 \frac{u^ke^{-(\sigma-\theta)u}}{k!}
 \le(\sigma-\theta)^{-k},
\]
which follows from the power series for the exponential.
Off this interval, use $|g|+g\le2C_0e^{2u}$ and
\eqref{eq:landau-tail-independent}. Since $\theta>0$, there is a constant
$C_2>0$, independent of $U,k,\eta,\gamma$, such that, for every
$\eta>0$ and arbitrarily large suitable $U$,
\begin{equation}\label{eq:landau-coeff-independent}
 |A_k(\sigma+i\gamma)|+\Rea A_k(\sigma)
 \le C_2\frac{e^{\eta U}}{(\sigma-\theta)^k}
 \quad\text{for all }\gamma\in\R,
 \qquad k\ge U.
\end{equation}

Suppose, for a contradiction, that a pole $\rho_0$ has
$\Rea\rho_0>\theta$. All poles have real part at most $2$, since the
original Laplace integral is holomorphic on $\Rea s>2$.
Let $Z$ be the pole set. Every compact subset of $\Rea s>0$ meets $Z$
in a finite set. If the closed disk with centre $\sigma$ and radius
$\sigma-\theta$ contains a pole, let $D$ be the distance to its nearest
pole; otherwise put $D=\sigma-\theta$. In either case
\[
 0<D\le\sigma-\theta,
 \qquad B(\sigma,D)\cap Z=\varnothing.
\]

First suppose there is a pole $\rho_1$ with
$\sigma-\Rea\rho_1<D$. There is then a real $\gamma$ for which
$z=\sigma+i\gamma$ has a unique nearest pole $\rho$ at distance $d<D$.
Here is the elementary justification. Start with
$\gamma=\Im\rho_1$ and vary it in a sufficiently small interval so
that $\rho_1$ remains at distance less than $D$. Only finitely many
poles can enter the disks of radius $D$ under this variation. For two
distinct such poles, equality of their distances from
$\sigma+i\gamma$ occurs at at most one value of $\gamma$: squared
distances have a difference affine in $\gamma$, and if their imaginary
parts coincide their different real parts, both less than $\sigma$,
give unequal distances for every $\gamma$. Avoid these finitely many
exceptional values. Choose $r$ with $d<r<D$ such that $B(z,r)$ contains
only $\rho$, and choose $r_1$ with $d<r_1<r$.

If $c\ne0$ is the residue at $\rho$, subtracting $c/(s-\rho)$ makes
$\Phi$ holomorphic on $B(z,r)$ after filling the removable singularity.
Cauchy's estimates consequently give constants $M,M'>0$ such that
\[
 \left|A_k(z)-\frac{c}{(z-\rho)^{k+1}}\right|\le Mr_1^{-k},
 \qquad |A_k(\sigma)|\le M'r_1^{-k}
 \quad(k\ge0).
\]
The second estimate uses the pole-free disk $B(\sigma,D)$.
Combining these inequalities with~\eqref{eq:landau-coeff-independent}
gives
\begin{equation}\label{eq:landau-final-independent}
 \frac{|c|}{d}\,d^{-k}
 \le (M+M')r_1^{-k}
       +C_2e^{\eta U}(\sigma-\theta)^{-k}.
\end{equation}

It remains to consider the case in which
$D\le\sigma-\Rea\rho$ for every pole $\rho$.
The pole-free closed-disk alternative is impossible because of
$\rho_0$. A nearest pole $\rho$ must satisfy
\[
 D\le\sigma-\Rea\rho\le|\sigma-\rho|=D.
\]
Hence it is the real pole $\beta=\sigma-D$, and it is the unique nearest
pole. Moreover $\beta>\theta$, since
$D\le\sigma-\Rea\rho_0<\sigma-\theta$.
Choose $D<r_1<r\le\sigma-\theta$ so that $B(\sigma,r)$ contains no
other pole. Its positive residue $c$ gives
\[
 \Rea A_k(\sigma)\ge cD^{-k-1}-Mr_1^{-k}
\]
by the same Cauchy estimate. Equation~\eqref{eq:landau-coeff-independent},
with $\gamma=0$, gives~\eqref{eq:landau-final-independent} with
$d=D$, $|c|=c$, and $M'=0$.

In either case $0<d<r_1\le\sigma-\theta$. Fix
$0<\eta<\log(r_1/d)$. Since $k\ge U$, multiplying
\eqref{eq:landau-final-independent} by $d^k$ shows
\[
 0<\frac{|c|}{d}
 \le (M+M'+C_2)
      \exp\bigl((\eta-\log(r_1/d))U\bigr).
\]
The right-hand side tends to zero along the arbitrarily large admissible
$U$, which is a contradiction.
\end{proof}


\begin{theorem}\label{thm:interval-grh}
Suppose that \eqref{eq:remaining-estimate} holds for every primitive
Dirichlet character, with $b>1$ fixed for each character. Then every
zero of every Dirichlet $L$-function in $0<\Re s<1$ has real part
$1/2$.
\end{theorem}

\begin{proof}
Fix $\chi$ and put $g(u)=H_\chi(e^u)$. Lemma~\ref{lem:positive-family}
supplies a meromorphic Laplace transform, with positive residues at
all real poles. Given $\eta>0$, choose
$0<\varepsilon<\eta/(2b)$. On any sufficiently large interval
$[Y,Y^b]$ in \eqref{eq:remaining-estimate}, set $U=\log Y$ and
$R_U(u)=e^{(1/2+\varepsilon)u}$. Then
\[
 \int_U^{bU}R_U(u)e^{-u/2}\,du
    =\frac{e^{b\varepsilon U}-e^{\varepsilon U}}{\varepsilon}
    \le \varepsilon^{-1}e^{b\varepsilon U}
    \le e^{\eta U}
\]
for all sufficiently large such $U$. Proposition~\ref{prop:landau-independent},
with $\theta=1/2$ and $\delta=b-1$, excludes every pole with
real part greater than $1/2$. Since $\chi\in\mathcal C_\chi$ and its
coefficient is positive, Lemma~\ref{lem:positive-family} excludes
every zero of $L(s,\chi)$ in that half-plane. The same argument
includes $\zeta$.

For a primitive character, a zero with $0<\Re s<1/2$ would give a
zero of $L(1-s,\bar\chi)$ with real part greater than $1/2$, by the
functional equation. The gamma factors are finite and nonzero
throughout the open critical strip. Thus every primitive zero there
has real part $1/2$.

Finally, if a character $\xi$ modulo $N$ is induced by $\xi^*$ of
conductor $f$, comparison of Euler factors, followed by analytic
continuation, gives
\[
 L(s,\xi)=L(s,\xi^*)
       \prod_{\substack{p\mid N\\p\nmid f}}
                    (1-\xi^*(p)p^{-s}).
\]
When $\Re s>0$, each factor is nonzero since
$|\xi^*(p)p^{-s}|=p^{-\Re s}<1$. Hence the result extends to all
characters. The usual Riemann hypothesis follows by taking conductor
one and using the classical location of the trivial zeta zeros.
\end{proof}

\begin{lemma}[Averaging nearby scales]\label{lem:logarithmic-averaging}
For $h>0$ define
\[
 \mathcal A_hH(x)=\frac1h\int_0^hH(xe^v)\,dv,
 \qquad \mathcal A_0H=H.
\]
If $g(u)=H(e^u)$ satisfies the transform and growth hypotheses of
Proposition~\ref{prop:landau-independent}, then so does
$g_h(u)=\mathcal A_hH(e^u)$. Every pole in $\Re s>0$ is preserved.
\end{lemma}

\begin{proof}
For $h>0$, Fubini's theorem initially in $\Re s>2$ gives
\[
 \Phi_h(s)=\frac{e^{hs}-1}{hs}\Phi(s)-E_h(s),\qquad
 E_h(s)=\frac1h\int_0^h e^{sv}\int_0^v g(w)e^{-sw}\,dw\,dv.
\]
The function $E_h$ is entire, since the integrations are over compact
intervals. The multiplier has removable value $1$ at zero, has no zero
in $\Re s>0$, and is positive on the positive real axis. It therefore
preserves all poles there, their simplicity, and positivity of real
residues. The bound $g_h(u)=O_h(e^{2u})$ follows directly. For $h=0$
there is nothing to prove.
\end{proof}

\begin{theorem}[Vanishing loss in an averaged estimate]
\label{thm:scale-averaged-criterion}
\label{thm:vanishing-loss-windows}
\label{cor:varying-positive-families}
Fix a primitive target character $\chi$, possibly of conductor one.
At each accuracy stage $j$, let
\[
 H_j(x)=\sum_{\nu\in\mathcal C_j}a_{j,\nu}\Re G_\nu(x),
 \qquad a_{j,\nu}>0,
\]
where $\mathcal C_j$ is a finite family of primitive characters containing
$\chi$ or $\bar\chi$.
Fix $b_j>1$, $h_j\ge0$ and $1/2<\theta_j<1$, with
$\theta_j\to1/2$. Put
\[
 I_j(Y)=\int_Y^{Y^{b_j}}(\mathcal A_{h_j}H_j(x))_+
                         x^{-\theta_j}\,\frac{dx}{x}.
\]
If for every $j$
\begin{equation}\label{eq:scale-averaged-target}
 \liminf_{Y\to\infty}\frac{\log(1+I_j(Y))}{\log Y}=0,
\end{equation}
then $L(s,\chi)$ and $L(s,\bar\chi)$ have no zeros to the right
of one half. In particular, it suffices to have a finite $M_j\ge0$
and arbitrarily late whole windows on which
$H_j(x)\le M_jx^{\theta_j}$.
\end{theorem}

\begin{proof}
Fix $j$. Lemma~\ref{lem:positive-family} supplies a meromorphic
Laplace transform for $H_j(e^u)$ in which every zero of the target
or its conjugate gives a nonzero residue. The preceding lemma
preserves these poles under the fixed local averaging.
For every $\eta>0$, \eqref{eq:scale-averaged-target} gives arbitrarily
large $U=\log Y$ such that the nonnegative majorant
$R_U(u)=(\mathcal A_{h_j}H_j(e^u))_+$ satisfies
\[
 \int_U^{b_jU}R_U(u)e^{-\theta_ju}\,du\le e^{\eta U}.
\]
Proposition~\ref{prop:landau-independent} excludes poles to the right
of $\theta_j$. Let $j$ tend to infinity. A fixed positive combination
is a special case, and the same reasoning then applies to every
character it contains. For the final pointwise sufficient condition,
choose $h_j=0$ and use $I_j(Y)\le M_j(b_j-1)\log Y$.
\end{proof}

\begin{corollary}[The combined signed arithmetic budget]
\label{cor:averaged-arithmetic-budget}
In the preceding theorem, let $\mathcal Y_j\subset(1,\infty)$ be
unbounded. Suppose a measurable, locally integrable real function $B_j$
satisfies $H_j\le B_j$ almost everywhere on
$[Y,e^{h_j}Y^{b_j}]$ for every $Y\in\mathcal Y_j$.
It suffices to establish \eqref{eq:scale-averaged-target}, with $Y$
restricted to this same set $\mathcal Y_j$ and with
\[
 \int_Y^{Y^{b_j}}(\mathcal A_{h_j}B_j(x))_+
                         x^{-\theta_j}\,\frac{dx}{x}
\]
in place of $I_j(Y)$. All signed contributions to $B_j$ may be
combined before the local average and the positive part are taken.
\end{corollary}

\begin{proof}
Positivity of the averaging operator gives
$(\mathcal A_{h_j}H_j)_+\le(\mathcal A_{h_j}B_j)_+$.
The expanded interval contains every argument used by that operator.
\end{proof}

\begin{remark}[Essential quantifiers and signs]
For each fixed accuracy index, the character family and its positive
coefficients must remain fixed while $Y$ tends to infinity through
whole admissible intervals. Replacing them at each $Y$ does not give
the repeated intervals required for a single Laplace transform.
The interval exponent and local averaging width must likewise be fixed
during each arbitrarily-large-$Y$ limit. A multiplicative constant in a
pointwise supplier must also be fixed during that limit.
Theorem~\ref{cor:varying-positive-families} allows all these data,
including the auxiliary characters and their coefficients, to vary
with the accuracy index, provided that the target character or its
conjugate retains a strictly positive coefficient. Only the total
exponent loss must tend to zero.

Two counterexamples delimit the analytic theorem. If
$g(u)=-e^{\beta u}$ with $\beta>\theta$, the majorant $R=0$ satisfies
every one-sided estimate but the transform has residue $-1$ at
$\beta$; thus positivity at real poles cannot be omitted.
If $g(u)=e^{\beta u}\cos u$, there are arbitrarily late negative
intervals of fixed length, although its transform has poles at
$\beta\pm i$. Such intervals cannot replace $[U,bU]$ with fixed
$b>1$. These are counterexamples to weaker criteria, not to the
criterion just proved.
\end{remark}

